\scanpage{47}\mspage{3}

% letterhead
{\footnotesize
\begin{tabular}{@{}l@{\quad}l@{}}
\deleted{MINISTERE DE L'EDUCATION NATIONALE} & \deleted{AVIGNON, le}\\
\deleted{ACADEMIE D'AIX-MARSEILLE} &\\
\hfil\rule{1.5cm}{0.4pt}\hfil &\\
\deleted{Inspection Académique} & \deleted{L'Inspecteur d'académie,}\\
\deleted{de Vaucluse} & \deleted{directeur des services départementaux}\\
\hfil\rule{1.5cm}{0.4pt}\hfil & \deleted{de l'Education Nationale de Vaucluse,}\\
\deleted{Tél. : 90 82 99 76} &\\
\deleted{Référence : Secrétariat n\textsuperscript{o}} & \deleted{à M}
\end{tabular}}
% end letterhead

donc $x$ n'est pas pt fixe pour $M$.

\deleted{On en conclut aussi que}
Le résultat $u_{E_0}=\mathrm{id}_{E_0}$ \inserted{$\forall u\in M$} s'écrit aussi
\[
\begin{aligned}
up(\unsure{x|n})&=p(x)\qquad\forall x\in E,\quad\text{i.e.}\\
up&=p\qquad\forall u\in M,
\end{aligned}
\]
ce qui prouve déjà que $M$ est cofiltrant
(car $u,v\in M$, on a $up=vp$, les deux mbres
étant égaux à $p$), \uline{qed}.

\deleted{J'ai envie de prouver ceci :}

\deleted{\uline{Proposition} 2 \inserted{(?)} Soit $X$ une catégorie \uline{1-connexe}
\uline{finie}, ayant un \deleted{objet} plus grand objet
$a$, et pseudo-cofiltrante. Alors $X$ est
cofiltrante.}

\deleted{On retrouve $(\unsure{\pm|+|\ddagger})$ la prop 1 en prenant
$X$ réduite à un seul objet. Je vais
essayer de calquer la dém. de prop 1.
Je note d'abord que, dans la catégorie
$X_{/a}$, \inserted{\unsure{quel que soit|quel que|qu'il est} \unsure{ds. cat.|cat.} $X$ finie,} il y a un plus petit objet,}

% archive sheet number: 46
